% !TEX root = ../main.tex
% !TEX program = lualatex
\documentclass[../main.tex]{subfiles}
\IfSubfilesClassLoaded{\externaldocument{\subfix{../build/main}}}{}
% =============================================================================
% CHAPTER 26: CURRENT EVENTS (2025)
% DESCRIPTION: Key 2025 updates on major acquisition and sustainment programs.
% =============================================================================
\begin{document}
\IfSubfilesClassLoaded{\chapter{EDO Study Guide}}{}

\chapter{Current Events in Navy Acquisition (2025)}\label{ch:current-events}
\addbibresource{assets/information/current_events.bib}

\section{Shipbuilding and Sustainment Updates}

\subsection{Columbia-Class (SSBN) Delays}
The Columbia-class ballistic missile submarine program, the Navy's highest acquisition priority, experienced significant schedule pressure in 2025. Navy officials reported an estimated 12- to 17-month delay in the delivery of the lead boat, USS \textit{District of Columbia} (SSBN-826), pushing its first planned patrol into 2030~\autocite{crs-columbia-2025}. The delays were attributed to a combination of factors, including shipyard workforce shortages, supply chain disruptions for critical components like the bow dome and main turbine generators, and the inherent complexities of a first-of-class build~\autocite{crs-columbia-2025,defensenews-columbia-2025}. By late 2025, General Dynamics Electric Boat reported the lead vessel was approximately 60\% complete, but the tight schedule margins remain a significant concern for maintaining the nation's continuous strategic deterrent capability~\autocite{usni-columbia-2025}.

\subsection{Amphibious Warship Block Buy (LPD-17)}
In a significant move to stabilize the amphibious ship industrial base, the Navy executed a multi-ship "block buy" procurement for four amphibious warships from HII Ingalls in late 2024 and early 2025~\autocite{breakingdefense-amphib-2024}. This procurement, valued at nearly \$1 billion in savings, includes three San Antonio-class amphibious transport docks (LPD-33, 34, and 35) and one America-class amphibious assault ship (LHA-10) to be built between FY25 and FY29~\autocite{breakingdefense-amphib-2024}. The block buy provides workload stability for the shipyard and its suppliers, enabling more efficient production and ensuring the Navy meets the legally mandated 31-ship amphibious force structure~\autocite{gao-amphib-2024}. This decision followed congressional pressure and signals a commitment to the amphibious fleet after a period of debate over future requirements.

\subsection{Shipyard Infrastructure Optimization Program (SIOP) Status}
The Shipyard Infrastructure Optimization Program (SIOP), a 20-year, \$21 billion effort to modernize the Navy's four public shipyards, reached a critical phase in 2025. A key milestone was the award of a \$377.7 million contract for seismic upgrades to Dry Dock \#4 at Puget Sound Naval Shipyard, which is essential for servicing nuclear submarines, including the future Columbia-class SSBNs~\autocite{kitsapsun-siop-2025}. However, the program faces significant cost growth due to inflation and unforeseen construction complexities, with annual construction cost escalation exceeding 6\%~\autocite{breakingdefense-siop-2025}. In July 2025, program officials announced a comprehensive review to revise cost and schedule estimates for the entire SIOP portfolio~\autocite{breakingdefense-siop-2025}. Despite these challenges, the Navy continues to invest approximately \$2.7 billion annually, with the goal of reducing submarine maintenance periods by up to 15\% and returning ships to the fleet faster~\autocite{breakingdefense-siop-2025}.

\end{document}
